\section{Protocol}
\label{sec:protocol}

\paragraph{Model and runtime.} Qwen3-8B \citep{yang2025qwen3} at a pinned revision
(Appendix~\ref{app:repro}), with thinking disabled through the chat template. Inference ran in float32
without quantization on one GPU. We chose float32 after finding that 8-bit and bfloat16 inference made
the log probability at a fixed position change by up to 0.843 and 0.875 nats when tokens were appended
after it. Float32 gave exactly zero (Appendix~\ref{app:precision}). Action events were scored with a
key--value-cache scorer. Before each run, the scorer was checked against one full forward pass per event
on identical prompts (maximum difference $5.08\times10^{-5}$ nats against a tolerance of 0.02). Answers
were generated greedily with at most 16 new tokens and parsed strictly: an answer counts only if, after
stripping outer whitespace, it is exactly one allowed label.

\paragraph{Splits and gate.} Three disjoint splits were generated from separate seeds and separate
entity-name pools. A development split (192 histories) was used only to check competence and answer
formats. A frozen gate split (384 histories) decided eligibility. The confirmation split (768 histories,
8{,}448 records, 11 per history) used a held-out paraphrase of the introduction and query. A task-cell
was \emph{competent} if strict direct retrieval of the current state was at least 0.95, current-only
decision accuracy at least 0.90, and invalid answers at most 0.05. Confirmation data were generated only
after the gate, recomputed from its saved scores, showed at least one competent cell. All cells are
reported. Task-level estimates average the fixed grid, and estimates from non-competent cells are labeled
as such.

\paragraph{Hypotheses and statistics.} The protocol fixed two primary hypotheses per task: H1,
$D_{\mathrm{superseded}} > 0$; and H2, $\Delta_{\mathrm{semantic}} \neq 0$, two-sided. Both use 97.5\%
percentile intervals from a history bootstrap (2{,}000 draws, fixed seed), Bonferroni-adjusted over the
two tasks. Secondary outcomes use unadjusted 95\% intervals: $\Delta_{\mathrm{mention}}$, the other
construction effects, net switch rates, and H3. Histories are the resampling unit. Both members, all
constructions and all queries of a history stay together. The two tasks are never pooled, and an
interval that includes zero is reported as inconclusive, not as evidence of no effect.

\paragraph{Prediction (H3).} Predictors that are available before the decision query is asked: the easy
query's margin between the current and the historical state, the historical candidate's log mass, easy
confidence, the edited token's position, prompt length, and design indicators. The comparison adds
$R_{\mathrm{easy}}$ to these. The outcome is a strict downstream failure on the historical constructions.
Models are L2-regularized logistic regressions ($C=1$, no tuning) evaluated by five-fold cross-validation
with folds grouped by history. Out-of-fold log loss with and without $R_{\mathrm{easy}}$ is compared by a
paired history bootstrap. The downstream answer's own confidence is reported separately as a
contemporaneous reference, not as a predictor, because it comes from the decision query itself.

\paragraph{Protocol history.} The confirmation protocol is the second redesign of a design that was
frozen before any model was run. The first version scored action events that the model never produces
and was corrected before inference. The second failed its competence gate because a trailing
\lit{Answer:} line made the model echo the prefix. The redesign changed only that line, the seeds and
the name pools. Appendix~\ref{app:history} lists every change and when it was made. The protocol records
are internal and dated. We call them \emph{frozen before confirmation}, not preregistered.
